\section{Configura\c{c}\~ao Experimental}
\label{sec:setup}

A avaliação da simulação mede se o perfil, os trechos recuperados e a CFG aumentam a capacidade do
modelo simulador de reconhecer o que a pessoa disse numa audiência que não entrou no seu perfil. A
fala gerada não serve como medida principal, porque o assunto da audiência, fornecido ao modelo, já
carrega a posição central de muitas opiniões: na audiência 1, do teste, o assunto registrado é
``Acusações de censura contra Alexandre de Moraes por exigir bloqueio de contas na rede social X'', e
uma das opiniões publicadas é ``Acusou Alexandre de Moraes de censura ao solicitar o bloqueio de
contas na rede social X''. Uma comparação entre a fala gerada e a opinião publicada daria crédito a
uma fala que apenas repetisse o assunto. A medida principal é, por isso, uma pergunta de múltipla
escolha, em que o assunto é o mesmo para todas as opções e o acerto ao acaso é conhecido (25\%).

\textbf{Perguntas.} Cada UDV de uma audiência de validação ou de teste cujo turno de evidência
pertence a um ator com perfil gera uma pergunta, ``Qual destas afirmações \textit{nome} fez nesta
audiência?'', acompanhada da data e do assunto da audiência e de quatro opções. A opção correta é a
opinião da UDV, na redação do resumo estruturado, e as demais são opiniões atribuídas pela matéria a
três outros participantes da mesma audiência, sorteados com semente fixa; opiniões da própria pessoa
e textos repetidos não entram, e a pergunta com menos de três outros participantes é descartada. Como
os distratores vêm da mesma audiência, eles tratam do mesmo tema, e a pergunta não pode ser resolvida
apenas pelo assunto. A UDV serve somente para ligar a opinião a um registro de ator pelo turno da
evidência, e a sentença escolhida como evidência não aparece na pergunta; a validade da pergunta
depende, portanto, da resolução da pessoa, e não da precisão da camada semântica. Das 106 opiniões de
teste ligadas a atores com perfil (Seção~\ref{sec:dados-particao}), 17 têm nível
\texttt{quote\_found}, 87 \texttt{semantic\_match\_high} e 2 \texttt{semantic\_match\_weak}.

\textbf{Leitura da resposta.} A escolha é lida na distribuição sobre o vocabulário da primeira
posição da resposta do modelo, restrita aos tokens das quatro letras, e antes da execução
confere-se que cada letra corresponde a um único token do tokenizador. Como LLMs favorecem certas
posições entre as opções~\citep{zheng2024large,pezeshkpour2024large}, cada pergunta é apresentada nas
quatro rotações cíclicas de uma ordem sorteada, de modo que cada opção ocupa uma vez cada letra. Em
cada rotação, as probabilidades das quatro letras são renormalizadas entre si, a probabilidade de
cada opção é a média das quatro rotações, e a resposta é a opção de maior probabilidade. As métricas
são o acerto e a probabilidade média atribuída à opção correta. Como essa leitura sempre escolhe uma
letra, mesmo quando o modelo iniciaria a resposta por outro token, e a probabilidade do primeiro
token pode divergir da resposta escrita~\citep{wang2024answer}, a soma das probabilidades das quatro
letras no vocabulário inteiro é registrada nas condições 0 a 2, como diagnóstico da parcela da
decisão do modelo que a leitura captura.

\textbf{Seleção de $k$ e $\gamma$.} Os dois hiperparâmetros são escolhidos apenas nas audiências de
validação, e antes da execução verifica-se que nenhum perfil e nenhuma fala usada na recuperação
contém audiência de validação ou de teste. Primeiro, $k \in \{3, 5, 10\}$ é o valor de maior acerto
da condição 2; em seguida, com esse $k$, $\gamma \in \{1;\ 1{,}5;\ 2;\ 3\}$ é o valor de maior acerto
da condição 3, para que a diferença entre as condições 2 e 3 dependa apenas da CFG. Empates são
decididos pela maior probabilidade média da opção correta e, persistindo, pelo menor valor da grade.
Como a grade inclui $\gamma = 1$, que reproduz a condição 2, a escolha desse valor indica que a CFG
não aumentou o acerto na validação. No teste, as condições são medidas apenas com o $k$ e o $\gamma$
escolhidos.

\textbf{Comparação entre condições.} Cada condição acrescenta um componente à anterior, e a
diferença de acerto entre condições vizinhas estima o efeito do componente acrescentado. As perguntas
de uma audiência compartilham assunto, falantes e distratores e, por isso, não são independentes; o
intervalo de 95\% de cada diferença é obtido por \textit{bootstrap} pareado de percentis, com 1.000
reamostragens de audiências inteiras~\citep{efron1993introduction}.

\textbf{Base da estimativa.} No teste, as condições 1 e 2 recebem também o pedido da base da
estimativa, sem as opções, para que a classificação não dependa dos distratores; a condição 3, que
tem o material da condição 2, herda a classificação dela. A classe é o primeiro rótulo que inicia a
resposta, e respostas sem rótulo são contadas à parte. Definiu-se antes da execução que as classes só
são informativas se o acerto na múltipla escolha diminuir de DIRETA para INDIRETA e de INDIRETA para
ESPECULATIVA.

\textbf{Geração aberta.} Para cada par formado por um ator com perfil e uma audiência de teste com
opinião ligada a ele, o modelo gera a base da estimativa, a fala e a justificação nas condições 1 a
3, com a data e o assunto da audiência. As falas não são comparadas às opiniões publicadas, pelo
motivo exposto no início desta seção; as medidas são o número de recusas, a fração de frases sem
fundamento na conferência da justificação e o número de respostas interrompidas pelo limite de
tokens (400 para a fala e 1.500 para a justificação). Para cada pedido gera-se também a fala com o
prompt negativo sozinho, lida para verificar se ele produz uma opinião genérica, como a CFG supõe, ou
uma recusa por falta de material, caso em que a CFG afastaria a fala da recusa, e não da opinião
padrão do modelo.

\textbf{Validação humana das UDVs.} A precisão da evidência das UDVs é medida numa amostra cega
sorteada só nas audiências de teste, porque as de validação foram usadas para comparar métodos. A
amostra tem 127 linhas (134 UDVs, 30 audiências), estratificadas em 35 das 41 UDVs de citação literal,
65 das 277 semânticas de alta confiança, 15 das 22 semânticas em que um prefixo curto da citação
ocorre na sentença escolhida, as 6 de baixa confiança e as 13 UDVs sem evidência, agrupadas em 6
linhas, uma por pessoa, nas quais o anotador indica se a pessoa falou na sessão. Nas demais linhas, o
anotador vê a opinião, a sentença e o texto ao redor, sem o nível, o cosseno ou o tipo de vínculo, e
julga a sentença como correta, parcial ou incorreta; ao menos 24 horas depois da primeira rodada, 20
dessas linhas voltam em outra ordem, para medir a concordância do anotador consigo mesmo. Os critérios foram registrados antes de qualquer julgamento: o limite inferior do intervalo
de Wilson~\citep{wilson1927probable} de 95\% deve atingir 0,90 para a precisão estrita (só
``correta'') da citação literal e 0,75 para a precisão tolerante (``correta'' ou ``parcial'') dos
casos semânticos de alta confiança.

\textbf{Recursos.} O modelo simulador é <modelo>, executado em <hardware>, e o codificador é o mesmo
Serafim 335m da UDV, na mesma versão de pesos. Cada execução registra o modelo, a versão dos prompts,
os identificadores dos tokens das letras e o SHA-256 dos arquivos de entrada.
